% ============================================================
%  6. MODELO Y GESTIÓN DE DATOS  (Subdocumento 5)
%  Informe 1 — Modelo de datos
%  Fuente: BTT Cap. 5 (RT-05.01 a RT-05.30)
% ============================================================
\chapter{Modelo y Gestión de Datos}
\label{sec:datos}

Este capítulo describe el modelo de datos de la solución, la selección de paradigmas de
persistencia, la gestión de calidad, migración, integración e interoperabilidad, y la
capa analítica, conforme al Capítulo 5 de las Bases Técnicas Transversales.

\section{Modelo de datos y diccionario}

La solución entregará el modelo de datos documentado y un diccionario de datos que
registre, para cada atributo: nombre, tipo, dominio de valores, obligatoriedad,
propietario y sensibilidad (RT-05.01). Los dominios de datos principales son:

\begin{itemize}
  \item \textbf{Transaccional:} pedidos, stock, trazabilidad de lote, temperatura,
        prueba de entrega, cobros y envases retornables.
  \item \textbf{Maestros:} productos, clientes, ubicaciones, conductores, vehículos,
        proveedores y rutas.
  \item \textbf{Analítico:} OTIF, \emph{fill rate} y costo de servir.
\end{itemize}

\section{Selección del paradigma de persistencia}

Se justifica la selección del paradigma y del motor de persistencia para cada dominio de
datos, con sus garantías transaccionales y su posición en el teorema CAP (RT-05.02):

\begin{table}[htbp]
\centering
\small
\caption{Paradigma y motor de persistencia por dominio}
\label{tab:persistencia}
\begin{tabularx}{\textwidth}{@{}L{4.4cm}L{4.4cm}X@{}}
\toprule
\textbf{Dominio} & \textbf{Paradigma / motor} & \textbf{Justificación CAP} \\
\midrule
Pedidos, stock y cobros &
Relacional ACID, alta disponibilidad multi-AZ &
Prioriza \emph{consistencia} y \emph{disponibilidad} en la partición (AP con réplica);
integridad transaccional crítica para no tener ``dos verdades''. \\
\addlinespace
Trazabilidad de lote y temperatura &
Relacional + series temporales (telemetría) &
Ingesta continua de sensores con consistencia eventual aceptable y alta disponibilidad;
reconstrucción auditada de la cadena de frío. \\
\addlinespace
Sincronización móvil offline &
SQLCipher local (AES-256) &
Consistencia local en el dispositivo con resolución determinista de conflictos por
\emph{timestamp} al sincronizar. \\
\addlinespace
Datos analíticos &
Almacén analítico separado &
Consultas analíticas nunca degradan la operación transaccional (RT-05.05). \\
\bottomrule
\end{tabularx}
\end{table}

\section{Trazabilidad, calidad de datos y seguridad}

\begin{itemize}
  \item \textbf{Trazabilidad de operaciones (RT-05.03):} toda operación de negocio será
        reconstruible --quién, qué, cuándo, desde qué dispositivo y con qué valores
        anteriores y posteriores-- durante el período de retención.
  \item \textbf{Calidad de datos (RT-05.04):} gestión conforme a ISO/IEC 25012, con
        validación en el punto de captura, indicadores de completitud, exactitud y
        consistencia, y tablero de calidad para el CLIENTE.
  \item \textbf{Exportación (RT-05.06):} exportación total de la información en formatos
        abiertos y documentados, sin costo y sin intervención del adjudicatario.
  \item \textbf{Retención y eliminación (RT-05.07):} política declarada por dominio de
        datos, coherente con la normativa (documentos tributarios 6 años, registros
        sanitarios 5 años, geolocalización 12 meses, Restricción R-41), con eliminación
        segura verificable.
  \item \textbf{Tratamiento de datos personales (RT-05.08):} seudonimización o cifrado a
        nivel de campo (AES-256) para categorías sensibles (comportamiento de pago,
        geolocalización de personas).
  \item \textbf{Gestión de datos maestros (RT-05.09):} estrategia que evita la duplicación
        de entidades compartidas (productos, clientes, conductores) entre módulos y con
        sistemas externos; prevalece el maestro interno de Puelche vía tabla de
        equivalencias (Decisión D11).
\end{itemize}

\section{Migración de datos}

Se presenta un plan de migración conforme a RT-05.11 a RT-05.15, con alcance, origen,
volumen, reglas de transformación, criterios de calidad, estrategia de ejecución y plan
de reversión. Incluye perfilado y saneamiento previo (RT-05.12), al menos dos ensayos
completos sobre preproducción (RT-05.13), conciliación cuantitativa y verificable
(RT-05.14), y acceso a los históricos no migrados durante la retención (RT-05.15). La
migración abarca los datos del WMS 2013, las planillas locales de Concepción y el ERP.

\section{Integración e interoperabilidad}

\begin{itemize}
  \item \textbf{Documentación de interfaces (RT-05.16):} servicios síncronos en OpenAPI
        3.1 y flujos dirigidos por eventos en AsyncAPI, generada desde el código.
  \item \textbf{Versionado y compatibilidad (RT-05.17):} contratos versionados
        semánticamente con compatibilidad hacia atrás y obsolescencia con preaviso de seis
        meses.
  \item \textbf{Autenticación entre sistemas (RT-05.18):} OAuth 2.1 con credenciales de
        cliente o autenticación mutua TLS; prohibida la clave estática en la URL.
  \item \textbf{Correlación (RT-05.19):} identificador de correlación común para seguir
        una operación a través de todos los sistemas.
  \item \textbf{Capa anticorrupción (RT-05.20):} aislamiento de integraciones con el ERP y
        el WMS 2013.
  \item \textbf{Declaración por integración (RT-05.21):} modo síncrono/asíncrono, volumen,
        ventana de disponibilidad y comportamiento ante no respuesta.
  \item \textbf{Intercambio masivo (RT-05.22):} carga y descarga masiva en formatos
        abiertos con validación, informe de errores y procesamiento parcial.
  \item \textbf{Estándares GS1 (RT-05.23):} GTIN, SSCC y GLN para interoperabilidad con
        el canal moderno y trazabilidad sanitaria.
\end{itemize}

\section{Analítica e inteligencia de negocio}

\begin{itemize}
  \item \textbf{Capa analítica (RT-05.25):} tableros operacionales y de gestión sobre los
        indicadores del caso (OTIF, \emph{fill rate}, costo de servir).
  \item \textbf{Exploración (RT-05.26):} filtros por período, unidad organizacional y
        dimensiones propias, con profundización del indicador agregado a la transacción.
  \item \textbf{Autoservicio (RT-05.27):} construcción de informes propios por el CLIENTE
        mediante herramienta de autoservicio con modelo semántico documentado.
  \item \textbf{Exportación y programación (RT-05.28):} informes exportables y
        programables por calendario.
  \item \textbf{Latencia (RT-05.29):} entre la transacción y su disponibilidad analítica
        no supera el umbral del caso (orientación: 4 horas).
  \item \textbf{Analítica predictiva (RT-05.30):} se valorará su incorporación con el
        modelo, variables, métrica de desempeño y plan de reentrenamiento documentados.
\end{itemize}

\vspace{4pt}
{\small\color{lafroxGrato}\pendiente{} Completar el diccionario de datos, el diagrama
entidad-relación, la estimación de volumetría (plantilla del Capítulo B de la BTT) y el
plan de migración detallado en la versión final.}
